\documentclass[conference,12pt]{IEEEtran}
\usepackage[utf8]{inputenc}
\usepackage[T1]{fontenc}
\usepackage[english]{babel}

\usepackage{amsmath}
\usepackage{graphicx}
\usepackage{enumitem}
\usepackage{todonotes}
\usepackage{subcaption}
\usepackage[hidelinks]{hyperref}

\begin{document}
    \title{Benchmarking AI Factories on MeluXina}
    \author{
        \IEEEauthorblockN{
            Jonah Holtmann,
            Mihkel Tiks,
            Luca Leonzio
        }
    }

    \maketitle

    \begin{abstract}

        We have created a microservice based AI Factory benchmarking suite that is able to efficiently capture and visualize key performance metrics to enable users to evaluate a given configurations performance on an HPC system. Through the use of Slurm and Docker we provide an easy to use framework that interfaces to commonly available tools and makes non-privileged installs possible. Further, the abstract SlurmServer class allows easy extendability for future tasks. The current set of benchmarks consists of IO500, ChromaDB, and vLLM.

    \end{abstract}

    \begin{IEEEkeywords}
        HPC, Benchmarking, IO500, Chroma, vLLM
    \end{IEEEkeywords}

    \section{Introduction}
    The goal of the EUMaster4HPC challenge is to create a benchmark suite deployable on MeluXina (and possibly other European HPC ressources) to evaluate the performance of the system.
    The students should become familiar with this emerging type of compute infrastructure on systems across the european union (EU).

    The Challenge is split into three phases: Onboarding, Prototyping, and Evaluation.
    Our teams' work is documented in a GitHub repository \cite{Team7}, which includes the codebase, documentation, and version control for our project.
    The end goal of the challenge is to provide a benchmark suite of the key AI factory components, including vector storage, LLM inference, and HPC performance.

    \section{Team Collaboration and Dynamics}

    The EUMaster4HPC organizers split the students into teams of four, with the goal of fostering collaboration and teamwork among students from different universities and backgrounds.
    While our team initially consisted of four members, one member stopped responding to our communications within the onboarding phase, leading to us continuing the challenge with three members.
    To ensure effective collaboration, we established a regular bi-weekly meeting schedule, which allowed us to balance the overhead of long meetings with the need of ensuring steady progress.
    The notes of each meeting are documented, so that team members can refer back as needed and also to ensure that members stay informed even if they miss a meeting.
    To divide the workload effectively, we assigned tasks and responsibilities based on each team member's strengths and interests.
    For this, each team member shared their background and experience, which helped us to identify the best fit for each area of the challenge.
    Further, we put an emphasis on ensuring that all team members equally contributed to the project design and implementation and that all members were involved in integration and testing.
    This ensures everyone has a good understanding of the entire benchmark suite and let us validate our work together.
    
    With the challenge consisting of three main areas of benchmarking (AI capabilities, vector storage, and HPC performance), we decided to split the responsibilities into these areas. 
    However, being assigned to a specific area did not mean that team members only worked on that area, but rather that they took ownership of it while also collaborating on the overall integration and testing of the benchmark suite.

    \subsection{Communication}
    To enable open and effective communication among team members, we established three platforms for communication and collaboration: a GitHub repository, a Discord group, and an instant messaging channel.

    The GitHub repository\cite{Team7} serves as the central hub for our codebase, documentation, and version control.
    We employ the GitHub Issue Tracker as an effective task management and progress tracking tool, which allows us to create and assign tasks, set deadlines, and track the progress of each task.
    Also, this ensures that all team members are aware of the current status of the project and can easily access the necessary information and resources.
    The issues are assigned according to the initial task distribution, but team members can also take on additional tasks as needed, which allows for flexibility and adaptability in our workflow.
    This also gives us the benefit of accountability and full transparency, with everyone knowing who has responsibility for which task and feature.
    The Discord group is used for group meetings and discussions.
    The video and screen sharing capabilities allow for easy collaboration and brainstorming sessions, which is especially important for the design and implementation of the benchmark suite.
    Finally, the instant messaging channel is utilized for quick questions and updates, ensuring that team members can stay connected and informed throughout the project.
    This approach came in handy especially during the onboarding and evaluation phases, where quick and easy communication and collaborative debugging is crucial to ensure swift progress.

    \subsection{Scheduling}

    As mentioned before, we established a regular bi-weekly meeting schedule to ensure effective collaboration and steady progress.
    Further, we set up a poll for possible meeting times, to make extraordinary meeting planning easier and help us to find the best timings to work together.
    Further, we went with a full commitment philosophy. Once a meeting time is set, we all commit, to make sure that nobody erroneously misses a meeting invitation.
    This let us reduce the overall amount of meetings, since there is minimal need to reschedule or repeat a meeting. Further, this also ensures that the meetings start without delay and that we could be productive right from the start of each meeting.


    \section{Tools and Software}

    \begin{figure*}[ht]
        \centering
        \includegraphics[width=.8\linewidth]{../docs/uml_diagram_simplified.pdf}
        \caption{Detail of the Architectural Design. The command-line interface (CLI) is the main entry point for users to interact with the benchmark suite. All workloads are started through it. A workload is defined by inheriting from the abstract SlurmServer class and allows easy extension to new workloads.}%
        \label{fig:uml_diagram}
    \end{figure*}

    
    The benchmark suite is designed to be modular and extensible, allowing for easy addition of new benchmarks and metrics as needed.
    It is based on a microservice architecture, with each benchmark and metric collection component running as a separate service.

    A simplified UML diagram of the architectural design is shown in Figure~\ref{fig:uml_diagram}..
    The user interacts with the benchmark suite through a command-line interface (CLI), which provides a simple and intuitive way to configure and run benchmarks, as well as to view results and metrics directly from the terminal.
    For this, the user needs a remote ssh connection since the CLI expects to be run on a node with direct access to the Slurm workload management.

    To make running the benchmarks as user-friendly as possible, we decided to use a combination of Slurm, Prometheus, Grafana, OpenTelemetry, and Docker.
    Since Slurm and Docker are usually available on HPC systems as standard tools, this ensures easy deployment and onboarding for users (and developers) of the benchmark suite.

    \subsection{Orchestration and Distribution}
    The main software tool for orchestration is Slurm, the popular workload management used on many HPC systems, including MeluXina.
    Slurm takes care of the workload distribution across the cluster, so there are no additional configuration steps to ensure benchmarks are not running with potentially invalid or duplicated resources.
    We interface to Slurm through the submission of normal Slurm job scripts.
    For our benchmarks, we supply relatively simple job scripts.
    
    Installation of the benchmark suite is handled through an installation script.
    This script expects only an internet connection and Docker to be available and should be run in advance of the first CLI startup. For file system benchmarking, a working compiler and a standard set of development tools are required, which may require the user to load some modules depending on their cluster configuration.
    Within the installation, the dependencies and used containers are fixed to to ensure a stable and consistent environment for the benchmarks.
    The installation script also takes care of setting up the Prometheus and Grafana containers, which are used for monitoring and visualization of the benchmark results.

    \subsection{Monitoring and Visualization}
    For the monitoring and data collection, we use Prometheus.
    This is flexible to use and allows us to easily collect and store the performance metrics of the benchmarks as well as the system metrics of the used nodes.
    Conceptually, once a job is started it makes its IP address known to the CLI which updates the Prometheus configuration so that the data is read.
    Also, we aggregate the metrics of all nodes attached to a benchmark job, so that there is less individual configuration and possibility for networking issues.
    Prometheus pulls these aggregations, so that there is no overhead introduced to a benchmark by the need of pushing some metrics over the network.

    For the visualization of the collected metrics, we use Grafana.
    Grafana lets us create custom dashboards to visualize the performance and system metrics in a clear and job specific way.
    This means that users can easily identify the job through its Slurm job ID and then view the metrics.
    For this purpose we use an exported Grafana dashboard, replacing the necessary variables and making it available to users through the CLI.

    Both Prometheus and Grafana are started through the CLI, and doing so toggles the hardware metric collection automatically.
    Once up and running, the user is given two ssh port forwarding commmands to be able to access Grafana and Prometheus from their local machine during the benchmarking process.
    Also available through the CLI is a command to export and save the collected data and logs to a compressed .zip archive, which can be used for later analysis and comparison of the benchmark results.

    \subsubsection{Hardware Metric Gathering}

    To gather the hardware metrics, we use a combination of a self built MPI aggregation service and OpenTelemetry to make job aggregates available to Prometheus while preserving all data for the time series data base.
    This lowers the overall overhead of hardware metric collection, but imposes a slight penalty on network heavy test scenarios.
    In our tested configurations, this accounts for less than 1\% deviation in performance metrics but may be more noticable with different configurations or workloads.
    By disabling metric gathering (also achieved by not starting the monitoring service) this overhead can be avoided.

    For custom scripts this tool can be started by adding a single line into the jobs Slurm script.

    

    \subsection{Benchmarking}
    As stated the benchmark suite is split into three main areas of benchmarking: AI inference, vector storage, and HPC performance (file system performance).
    To make the benchmark suite extendable and modular, we decided to implement an abstract SlurmServer class, which defines the default methods, properties and behaviour of a benchmark job (refer to Figure~\ref{fig:uml_diagram}).

    For this purpose, the SlurmServer class defines four variables: the Slurm job ID, the ip address of the master node of the job, a flag indicating job status, and a flag indicating if the setup is complete and benchmarking can start.
    This allows the CLI to easily track the status of a job and allows us to act as a soft switch to prohibit a benchmark being started before the necessary setup is complete.
    The SlurmServer also provides the CLI with convenient method hooks to start and safely stop a benchmark job, both of which can be configured to be blocking or non-blocking, and also to overwrite or extend the default behaviour.
    This allows us to handle cases where in addition to a workload there also exists a service at which it is directed, which may be started together with the routine.

    Further, the SlurmServer provides a method to check a status to let the user manually check for a services health. This uses an internal readiness check which needs to be overwritten by the specific benchmark implementation, to ensure correct setup of the two indicator flags.



    \subsubsection{AI Inference}
    The AI inference benchmark uses vLLM, a high-performance inference engine for large language models (LLMs), to evaluate the performance of LLM inference on MeluXina.
    By default, the benchmark runs a predefined set of inference tasks, which can be easily extended by users to include custom tasks and models.
    To pull a model the user needs to initialize HF\_TOKEN in their environment, which is used to pull the corresponding model from Hugging Face.
    Therefor, the benchmark triggers the two-step launch process of SlurmServer, indicating running once the Slurm job allocation is launched and then indicating setup complete once the vLLM server is up and running and ready to receive inference requests.
    During the start on the CLI the user can specify how many nodes to use for the benchmark, specify a distributed server list and specify the model to be used for inference.
    Throguh the wide variety of configuration settings, users can easily test and compare the performance of different models and configurations on MeluXina, which can provide valuable insights into the capabilities and limitations of the system for AI workloads.
    From our tests, the vLLM benchmark scales well to many nodes and inference servers.
    
    To interface with Prometheus, we use vLLM's built-in metrics endpoint, only pointing Prometheus to the master node.

    \subsubsection{Vector Storage}
    The vector storage benchmark uses ChromaDB, a popular vector storage data base, to evaluate MeluXina's capabilities.
    Similar to vLLM it sets up a ChromaDB server through docker and--once started--signals to the CLI that benchmarking can commence.
    It offers a wide variety of telemetrics through its own metrics endpoint, which is scraped by Prometheus and made available as a Grafana Dashboard.
    The user can configure the ressources and setup for the benchmarking run.
    The benchmark scales well to many nodes, enabling the user to accurately study MeluXina's performance for their given workload size.

    \subsubsection{File System Performance}
    The file system performance benchmark uses IO500, a widely used benchmark suite for evaluating the performance of parallel file systems in HPC environments, to evaluate the performance of MeluXina's file system.
    MeluXina uses LustreFS for its main work storage pool and allows users to configure a directory in Lustre.
    The IO500 benchmark suite is installed in the initial setup of our benchmark suite.
    By default, a predefined set of IO500 benchmarks is run, which has been taken from the IO500 challenge of the Student Cluster Competitions at the ISC HPC series~\cite{ischpc}.

    This configuration is designed to run within a reasonable time frame but may require some hundreds of GB of storage space depending on the users file system speed, so it should not be run in a quota constrained setting. Doing so may yield unrecoverable errors, leading to aborting the benchmark run.

    During the launch of the benchmark, the user is asked for a directory to run the benchmark in.
    This directory is monitored for the benchmark job and will be written to, read from and possibly cleaned by the benchmark, so caution is advised. By default it is set to a subdirectory of MeluXina's scratch directory.
    The user is then asked for their own IO500 configuration file, allowing easy customization to test specific scenarios as needed.
    While the benchmark is running, the user can only see the node level metrics, since the IO500 benchmark does not provide a fine-grained metric during the run, but only a final report at the end of the benchmark. Once the benchmark is finished, the user is shown the final report in the CLI, which also allows to select previous runs and compare performance.
    On benchmarking end, the user is notified through the CLI's check status method, which displays the progress and final report of (all) ongoing IO500 benchmark runs.

    \section{Key Takeaways and Lessons Learned}

    Working on this challenge has been a valuable learning experience for our team, providing insights into the complexities of benchmarking AI factories on HPC systems and the importance of effective collaboration and communication in a team setting.
    Working with the extensive and multi-faceted soft- and hardware infrastructure of MeluXina has been a great opportunity to gain hands-on experience.
    Being able to apply techniques and tools we have learned during our classes to a system has helped to solidify our understanding and gain a feeling for the practical implications of performance engineering and benchmarking.
    Getting to run large node allocations has taught us how to efficiently handle compute ressources, so that we do not block large partitions for setup tasks.

    Through working as a team each member has been able to both contribute their own strengths and also step out of their comfort zone to learn new skills and gain experience in areas they were less familiar with, which has been a great opportunity for personal growth and development.
    A key takeaway is the importance of clear and effective communication to avoid errors previously encountered by others and to ensure that all team members are on the same page regarding the project goals, tasks, and progress.

    Being proactive and committing to meetings is important to react to team members becoming unresponsive and avoids work being left incomplete.
    

\bibliographystyle{IEEEtran}
\bibliography{references}

\end{document}